\documentclass[11pt]{article}
\usepackage[margin=1in]{geometry}
\usepackage{amsmath,amssymb,amsthm,mathtools,microtype}
\usepackage[hidelinks]{hyperref}

\newtheorem{theorem}{Theorem}
\newtheorem{lemma}{Lemma}
\newcommand{\C}{\mathbb C}
\newcommand{\N}{\mathbb N}

\title{Riesz-complete extensions realize every Weyl order}
\author{Candidate full solution to an open question in arXiv:2510.20249}
\date{17 August 2026}

\begin{document}
\maketitle

\begin{center}
\fbox{\parbox{0.92\textwidth}{\textbf{Result.}
The existence of a Riesz-complete non-self-adjoint extension imposes no
restriction on Weyl order.  For every $\rho\in[0,\infty]$ there is an entire
simple symmetric operator with deficiency indices $(1,1)$, of Weyl order
$\rho$, which has such an extension.}}
\end{center}

\section{The source question and its criterion}

For an entire simple symmetric operator $T$ with deficiency indices $(1,1)$,
the source defines the Weyl height $h_T(r)$ and Weyl order
\[
 \rho_T=\limsup_{r\to\infty}\frac{\log h_T(r)}{\log r}.
\]
At the end of its section on completeness (source PDF page 45), it asks
whether the existence of a Riesz-complete non-self-adjoint extension restricts
$\rho_T$.

We use two results proved in the source \cite{Wang}.  If the eigenvalues of a
strictly accumulative extension $T_\infty$ are
$\lambda_n\in\C_-$ and the contractive Weyl function is the corresponding
meromorphic Blaschke product, then
\begin{equation}\label{eq:rieszcriterion}
 T_\infty\text{ is Riesz complete}
 \quad\Longleftrightarrow\quad
 \inf_n\prod_{m\ne n}
 \left|\frac{\overline{\lambda_n}-\overline{\lambda_m}}
 {\overline{\lambda_n}-\lambda_m}\right|>0.
\end{equation}
Moreover, writing $B$ for the contractive Weyl function,
\begin{equation}\label{eq:height}
 h_T(r)=\frac{1}{2\pi i}\int_0^r\frac{dt}{t}
       \int_{-t}^{t}d\log B(u)+O(1).
\end{equation}
The criterion is Theorem 6.12 on source PDF page 44 and the height formula is
on source PDF page 15.  The source also establishes on PDF page 12 that the
corresponding entire Nevanlinna curves are realized by entire operators of
this kind.  Thus it is
enough to construct a Blaschke product satisfying \eqref{eq:rieszcriterion}
and then compute the order in \eqref{eq:height}.

\section{Idea of the proof}

Choose the real parts $x_n$ to have whatever radial density is needed for the
desired order.  This could destroy interpolation if the $x_n$ become crowded.
The remedy is to put the poles only an exponentially tiny distance below the
real line: if $d_n$ is the nearest-neighbor gap at $x_n$, use
$y_n=2^{-n}d_n$.  For two conjugate poles, the failure of their
pseudohyperbolic distance to equal one is then at most $4\,2^{-n-m}$.
Summing logarithms gives a uniform positive lower bound for the full
separation product, independently of how dense the real parts become.

Each pole still contributes a fixed positive amount to the phase integral
once the integration interval passes its real part.  Conversely, its
contribution before that point is controlled by $y_n/x_n$.  Hence the Weyl
height has the same power-order as the counting density of the $x_n$.  Three
explicit choices of $x_n$ cover positive finite, zero, and infinite order.

\section{Uniform interpolation with arbitrary radial density}

Let $0<x_1<x_2<\cdots$ tend to infinity and set
\[
 d_n=\min\left\{1,\inf_{m\ne n}|x_n-x_m|\right\},\qquad
 y_n=2^{-n}d_n,\qquad \lambda_n=x_n-i y_n.
\]
Since the sequence is discrete, $d_n>0$.  Put $z_n=\overline{\lambda_n}$ and
define
\begin{equation}\label{eq:B}
 B(\lambda)=\prod_{n=1}^{\infty}
 \frac{\lambda-z_n}{\lambda-\overline{z_n}}.
\end{equation}
The harmless usual unimodular normalizing factors may be inserted.  Because
$\sum_n y_n<\infty$ and $x_n\to\infty$, the product converges normally away
from its poles.  It is a meromorphic inner function on the upper half-plane,
with zeros $z_n$, poles $\lambda_n$, and zero exponential factor.

\begin{lemma}[Uniform separation]\label{lem:separation}
The sequence $(z_n)$ satisfies
\[
 \inf_n\prod_{m\ne n}
 \left|\frac{z_n-z_m}{z_n-\overline{z_m}}\right|\ge e^{-2}.
\]
Consequently the extension with eigenvalues $(\lambda_n)$ is Riesz complete.
\end{lemma}

\begin{proof}
For $m\ne n$, write
\[
 r_{nm}=\left|\frac{z_n-z_m}{z_n-\overline{z_m}}\right|.
\]
A direct calculation gives
\[
 r_{nm}^2=1-u_{nm},\qquad
 u_{nm}=\frac{4y_ny_m}
 {(x_n-x_m)^2+(y_n+y_m)^2}.
\]
The definition of $d_n$ gives
$|x_n-x_m|\ge\max\{d_n,d_m\}$ and therefore
\[
 0\le u_{nm}\le4\,2^{-n-m}.
\]
Because $n,m\ge1$ and $n\ne m$, the right side is at most $1/2$.  The
elementary inequality $-\frac12\log(1-u)\le u$ for $0\le u\le1/2$ now yields
\[
 -\log\prod_{m\ne n}r_{nm}
 \le\sum_{m\ne n}u_{nm}
 <4\,2^{-n}\sum_{m=1}^{\infty}2^{-m}
 =4\,2^{-n}\le2.
\]
This proves the product bound.  The Riesz-completeness conclusion is exactly
the source criterion \eqref{eq:rieszcriterion}.
\end{proof}

\section{The Weyl-height estimate}

For real $u$, logarithmically differentiating \eqref{eq:B} gives a positive
phase density,
\[
 \frac1i\,d\log B(u)
 =\sum_{n=1}^{\infty}\frac{2y_n}{(u-x_n)^2+y_n^2}\,du.
\]
Define
\[
 a_n(t)=\frac1{2\pi}\int_{-t}^{t}
          \frac{2y_n}{(u-x_n)^2+y_n^2}\,du,
 \qquad 0\le a_n(t)\le1.
\]
Tonelli's theorem and \eqref{eq:height} give
\begin{equation}\label{eq:H}
 h_T(r)=\sum_n\int_0^r a_n(t)\frac{dt}{t}+O(1).
\end{equation}

\begin{lemma}[Density controls height]\label{lem:height}
There are absolute constants $c,C>0$ such that, for all sufficiently large
$r$,
\begin{align}
 h_T(r)&\ge c\,\#\{n:x_n\le r/3\}-C,\label{eq:lower}\\
 h_T(r)&\le
 \sum_{x_n\le2r}\log\frac{2r}{x_n}+C.\label{eq:upper}
\end{align}
\end{lemma}

\begin{proof}
If $x_n\le r/3$, then $x_n+y_n\le r/2$ for large $r$.  For every
$t\in[r/2,r]$, the interval $[-t,t]$ contains
$[x_n-y_n,x_n+y_n]$.  Integration on this smaller interval shows
$a_n(t)\ge1/2$.  Its contribution to \eqref{eq:H} is therefore at least
$\frac12\log2$, proving \eqref{eq:lower}.

For the upper bound, if $t\le x_n/2$, then $|u-x_n|\ge x_n/2$ on
$[-t,t]$, whence
\[
 a_n(t)\le \frac{8}{\pi}\frac{t y_n}{x_n^2}.
\]
For $t>x_n/2$ use only $a_n(t)\le1$.  The contribution of an index with
$x_n\le2r$ is at most
$C y_n/x_n+\log(2r/x_n)$; the first terms are summable.  If $x_n>2r$, its
contribution is at most $Cr y_n/x_n^2$, and their sum is bounded (indeed it is
$O(r^{-1})$).  Substitution in \eqref{eq:H} proves \eqref{eq:upper}.
\end{proof}

\section{Complete realization of all Weyl orders}

\begin{theorem}[Full answer to the source question]\label{thm:main}
For every $\rho\in[0,\infty]$ there is a transcendental entire simple
symmetric operator $T$ with deficiency indices $(1,1)$ and Weyl order
$\rho_T=\rho$ which has a Riesz-complete non-self-adjoint extension.
\end{theorem}

\begin{proof}
Use the construction above.  Lemma \ref{lem:separation} always gives the
Riesz-complete extension $T_\infty$.  Its eigenvalues lie strictly in
$\C_-$, so it is non-self-adjoint.

For $0<\rho<\infty$, take $x_n=n^{1/\rho}$.  The counting term in
\eqref{eq:lower} is comparable to $r^\rho$, while
\[
 \sum_{n\le(2r)^\rho}\log\frac{2r}{n^{1/\rho}}=O(r^\rho)
\]
by Stirling's formula.  Lemma \ref{lem:height} therefore gives
$c_\rho r^\rho\le h_T(r)\le C_\rho r^\rho$ for large $r$, and hence
$\rho_T=\rho$.

For $\rho=0$, take $x_n=e^{\sqrt n}$.  Then
\[
 c(\log r)^2\le h_T(r)\le C(\log r)^3.
\]
Thus $h_T(r)/\log r\to\infty$, so the operator is transcendental, while
$\log h_T(r)/\log r\to0$; hence $\rho_T=0$.

Finally, for $\rho=\infty$, take $x_n=\log(n+1)$.  The lower bound gives
$h_T(r)\ge c e^{r/3}$ for large $r$, and consequently
$\log h_T(r)/\log r\to\infty$.  This is infinite Weyl order.

All possible values of Weyl order have now been realized, so Riesz
completeness imposes no restriction on it.
\end{proof}

\section{Verification and novelty scope}

The included verifier evaluates finite separation products for all three
growth regimes and checks the analytic bound $u_{nm}\le4\,2^{-n-m}$.  It also
computes truncated phase integrals from \eqref{eq:H}; these checks are not
used in the proof.

A bounded search through 17 August 2026 used the exact paper title, the exact
question sentence, and combinations of ``Riesz complete'', ``Weyl order'',
and ``interpolating sequence''.  It found no later answer; arXiv still lists
only version 1 of the source.  The result uses the source's realization,
Riesz-basis criterion, and height formula.  The interpolating construction and
the all-orders conclusion are proved here.

\begin{thebibliography}{9}
\bibitem{Wang}
Y.~Wang,
\emph{Complex analysis of symmetric operators. II: entire operators with
deficiency index 1}, arXiv:2510.20249v1 (2025).
\end{thebibliography}

\end{document}
